%% bare_conf.tex
%% V1.4b
%% 2015/08/26
%% by Michael Shell
%% See:
%% http://www.michaelshell.org/
%% for current contact information.
%%
%% This is a skeleton file demonstrating the use of IEEEtran.cls
%% (requires IEEEtran.cls version 1.8b or later) with an IEEE
%% conference paper.
%%
%% Support sites:
%% http://www.michaelshell.org/tex/ieeetran/
%% http://www.ctan.org/pkg/ieeetran
%% and
%% http://www.ieee.org/

%%*************************************************************************
%% Legal Notice:
%% This code is offered as-is without any warranty either expressed or
%% implied; without even the implied warranty of MERCHANTABILITY or
%% FITNESS FOR A PARTICULAR PURPOSE! 
%% User assumes all risk.
%% In no event shall the IEEE or any contributor to this code be liable for
%% any damages or losses, including, but not limited to, incidental,
%% consequential, or any other damages, resulting from the use or misuse
%% of any information contained here.
%%
%% All comments are the opinions of their respective authors and are not
%% necessarily endorsed by the IEEE.
%%
%% This work is distributed under the LaTeX Project Public License (LPPL)
%% ( http://www.latex-project.org/ ) version 1.3, and may be freely used,
%% distributed and modified. A copy of the LPPL, version 1.3, is included
%% in the base LaTeX documentation of all distributions of LaTeX released
%% 2003/12/01 or later.
%% Retain all contribution notices and credits.
%% ** Modified files should be clearly indicated as such, including  **
%% ** renaming them and changing author support contact information. **
%%*************************************************************************

% *** Authors should verify (and, if needed, correct) their LaTeX system  ***
% *** with the testflow diagnostic prior to trusting their LaTeX platform ***
% *** with production work. The IEEE's font choices and paper sizes can   ***
% *** trigger bugs that do not appear when using other class files.       ***                          ***
% The testflow support page is at:
% http://www.michaelshell.org/tex/testflow/

\documentclass[conference, compsoc, twoside]{IEEEtran}
% Some Computer Society conferences also require the compsoc mode option,
% but others use the standard conference format.
%
% If IEEEtran.cls has not been installed into the LaTeX system files,
% manually specify the path to it like:
% \documentclass[conference]{../sty/IEEEtran}

% Some very useful LaTeX packages include:
% (uncomment the ones you want to load)


% *** MISC UTILITY PACKAGES ***
%
%\usepackage{ifpdf}
% Heiko Oberdiek's ifpdf.sty is very useful if you need conditional
% compilation based on whether the output is pdf or dvi.
% usage:
% \ifpdf
%   % pdf code
% \else
%   % dvi code
% \fi
% The latest version of ifpdf.sty can be obtained from:
% http://www.ctan.org/pkg/ifpdf
% Also, note that IEEEtran.cls V1.7 and later provides a builtin
% \ifCLASSINFOpdf conditional that works the same way.
% When switching from latex to pdflatex and vice-versa, the compiler may
% have to be run twice to clear warning/error messages.

% *** CITATION PACKAGES ***
%
%\usepackage{cite}
%\usepackage[noadjust]{cite}
\usepackage[numbers,sort&compress]{natbib}
\usepackage[labelsep=period,labelfont=bf,justification=justified,font=small]{caption}
\usepackage{balance}
\usepackage{graphicx}
%\bibliographystyle{ieeetr}
%\usepackage[sorting=none]{biblatex}
%\bibliography{journals,phd-references}
% cite.sty was written by Donald Arseneau
% V1.6 and later of IEEEtran pre-defines the format of the cite.sty package
% \cite{} output to follow that of the IEEE. Loading the cite package will
% result in citation numbers being automatically sorted and properly
% "compressed/ranged". e.g., [1], [9], [2], [7], [5], [6] without using
% cite.sty will become [1], [2], [5]--[7], [9] using cite.sty. cite.sty's
% \cite will automatically add leading space, if needed. Use cite.sty's
% noadjust option (cite.sty V3.8 and later) if you want to turn this off
% such as if a citation ever needs to be enclosed in parenthesis.
% cite.sty is already installed on most LaTeX systems. Be sure and use
% version 5.0 (2009-03-20) and later if using hyperref.sty.
% The latest version can be obtained at:
% http://www.ctan.org/pkg/cite
% The documentation is contained in the cite.sty file itself.

% *** GRAPHICS RELATED PACKAGES ***
%
\ifCLASSINFOpdf
  % \usepackage[pdftex]{graphicx}
  % declare the path(s) where your graphic files are
  %\graphicspath{{pdf/}{jpeg/}}
  % and their extensions so you won't have to specify these with
  % every instance of \includegraphics
  % \DeclareGraphicsExtensions{.pdf,.jpeg,.png}
  %\DeclareGraphicsExtensions{.pdf,.jpeg,.png}
\else
  % or other class option (dvipsone, dvipdf, if not using dvips). graphicx
  % will default to the driver specified in the system graphics.cfg if no
  % driver is specified.
  % \usepackage[dvips]{graphicx}
  % declare the path(s) where your graphic files are
  % \graphicspath{{../eps/}}
  % and their extensions so you won't have to specify these with
  % every instance of \includegraphics
  % \DeclareGraphicsExtensions{.eps}
  %\usepackage{subfigure}
\fi
% graphicx was written by David Carlisle and Sebastian Rahtz. It is
% required if you want graphics, photos, etc. graphicx.sty is already
% installed on most LaTeX systems. The latest version and documentation
% can be obtained at: 
% http://www.ctan.org/pkg/graphicx
% Another good source of documentation is "Using Imported Graphics in
% LaTeX2e" by Keith Reckdahl which can be found at:
% http://www.ctan.org/pkg/epslatex
%
% latex, and pdflatex in dvi mode, support graphics in encapsulated
% postscript (.eps) format. pdflatex in pdf mode supports graphics
% in .pdf, .jpeg, .png and .mps (metapost) formats. Users should ensure
% that all non-photo figures use a vector format (.eps, .pdf, .mps) and
% not a bitmapped formats (.jpeg, .png). The IEEE frowns on bitmapped formats
% which can result in "jaggedy"/blurry rendering of lines and letters as
% well as large increases in file sizes.
%
% You can find documentation about the pdfTeX application at:
% http://www.tug.org/applications/pdftex



% *** MATH PACKAGES ***
%
\usepackage{amsmath}
\DeclareMathOperator*{\argmin}{arg\,min} % thin space, limits underneath in displays
\DeclareMathOperator*{\argmax}{arg\,max} % thin space, limits underneath in displays
\newcommand{\ineq}{%
	\mathrel{\mkern1mu\underline{\mkern-1mu\in\mkern-1mu}\mkern1mu}}
%\DeclareMathOperator*{\argmin}{argmin} % no space, limits underneath in displays
%\DeclareMathOperator{\argmin}{arg\,min} % thin space, limits on side in displays
%\DeclareMathOperator{\argmin}{argmin} % no space, limits on side in displays
% A popular package from the American Mathematical Society that provides
% many useful and powerful commands for dealing with mathematics.
%
% Note that the amsmath package sets \interdisplaylinepenalty to 10000
% thus preventing page breaks from occurring within multiline equations. Use:
%\interdisplaylinepenalty=2500
% after loading amsmath to restore such page breaks as IEEEtran.cls normally
% does. amsmath.sty is already installed on most LaTeX systems. The latest
% version and documentation can be obtained at:
% http://www.ctan.org/pkg/amsmath

% *** SPECIALIZED LIST PACKAGES ***
%
% \usepackage{algorithmic}
% algorithmic.sty was written by Peter Williams and Rogerio Brito.
% This package provides an algorithmic environment fo describing algorithms.
% You can use the algorithmic environment in-text or within a figure
% environment to provide for a floating algorithm. Do NOT use the algorithm
% floating environment provided by algorithm.sty (by the same authors) or
% algorithm2e.sty (by Christophe Fiorio) as the IEEE does not use dedicated
% algorithm float types and packages that provide these will not provide
% correct IEEE style captions. The latest version and documentation of
% algorithmic.sty can be obtained at:
% http://www.ctan.org/pkg/algorithms
% Also of interest may be the (relatively newer and more customizable)
% algorithmicx.sty package by Szasz Janos:
% http://www.ctan.org/pkg/algorithmicx

% *** ALIGNMENT PACKAGES ***
%
\usepackage{array}
% Frank Mittelbach's and David Carlisle's array.sty patches and improves
% the standard LaTeX2e array and tabular environments to provide better
% appearance and additional user controls. As the default LaTeX2e table
% generation code is lacking to the point of almost being broken with
% respect to the quality of the end results, all users are strongly
% advised to use an enhanced (at the very least that provided by array.sty)
% set of table tools. array.sty is already installed on most systems. The
% latest version and documentation can be obtained at:
% http://www.ctan.org/pkg/array

% IEEEtran contains the IEEEeqnarray family of commands that can be used to
% generate multiline equations as well as matrices, tables, etc., of high
% quality.

% *** SUBFIGURE PACKAGES ***
%\ifCLASSOPTIONcompsoc
%  \usepackage[caption=false,font=normalsize,labelfont=sf,textfont=sf]{subfig}
%\else
%  \usepackage[caption=false,font=footnotesize]{subfig}
%\fi
% subfig.sty, written by Steven Douglas Cochran, is the modern replacement
% for subfigure.sty, the latter of which is no longer maintained and is
% incompatible with some LaTeX packages including fixltx2e. However,
% subfig.sty requires and automatically loads Axel Sommerfeldt's caption.sty
% which will override IEEEtran.cls' handling of captions and this will result
% in non-IEEE style figure/table captions. To prevent this problem, be sure
% and invoke subfig.sty's "caption=false" package option (available since
% subfig.sty version 1.3, 2005/06/28) as this is will preserve IEEEtran.cls
% handling of captions.
% Note that the Computer Society format requires a larger sans serif font
% than the serif footnote size font used in traditional IEEE formatting
% and thus the need to invoke different subfig.sty package options depending
% on whether compsoc mode has been enabled.
%
% The latest version and documentation of subfig.sty can be obtained at:
% http://www.ctan.org/pkg/subfig

% *** FLOAT PACKAGES ***
%
%\usepackage{fixltx2e}
% fixltx2e, the successor to the earlier fix2col.sty, was written by
% Frank Mittelbach and David Carlisle. This package corrects a few problems
% in the LaTeX2e kernel, the most notable of which is that in current
% LaTeX2e releases, the ordering of single and double column floats is not
% guaranteed to be preserved. Thus, an unpatched LaTeX2e can allow a
% single column figure to be placed prior to an earlier double column
% figure.
% Be aware that LaTeX2e kernels dated 2015 and later have fixltx2e.sty's
% corrections already built into the system in which case a warning will
% be issued if an attempt is made to load fixltx2e.sty as it is no longer
% needed.
% The latest version and documentation can be found at:
% http://www.ctan.org/pkg/fixltx2e

%\usepackage{stfloats}
% stfloats.sty was written by Sigitas Tolusis. This package gives LaTeX2e
% the ability to do double column floats at the bottom of the page as well
% as the top. (e.g., "\begin{figure*}[!b]" is not normally possible in
% LaTeX2e). It also provides a command:
%\fnbelowfloat
% to enable the placement of footnotes below bottom floats (the standard
% LaTeX2e kernel puts them above bottom floats). This is an invasive package
% which rewrites many portions of the LaTeX2e float routines. It may not work
% with other packages that modify the LaTeX2e float routines. The latest
% version and documentation can be obtained at:
% http://www.ctan.org/pkg/stfloats
% Do not use the stfloats baselinefloat ability as the IEEE does not allow
% \baselineskip to stretch. Authors submitting work to the IEEE should note
% that the IEEE rarely uses double column equations and that authors should try
% to avoid such use. Do not be tempted to use the cuted.sty or midfloat.sty
% packages (also by Sigitas Tolusis) as the IEEE does not format its papers in
% such ways.
% Do not attempt to use stfloats with fixltx2e as they are incompatible.
% Instead, use Morten Hogholm'a dblfloatfix which combines the features
% of both fixltx2e and stfloats:
%
% \usepackage{dblfloatfix}
% The latest version can be found at:
% http://www.ctan.org/pkg/dblfloatfix

\usepackage[english]{babel}
\usepackage{blindtext}

% *** PDF, URL AND HYPERLINK PACKAGES ***
%
\usepackage{url}
% url.sty was written by Donald Arseneau. It provides better support for
% handling and breaking URLs. url.sty is already installed on most LaTeX
% systems. The latest version and documentation can be obtained at:
% http://www.ctan.org/pkg/url
% Basically, \url{my_url_here}.

% *** Do not adjust lengths that control margins, column widths, etc. ***
% *** Do not use packages that alter fonts (such as pslatex).         ***
% There should be no need to do such things with IEEEtran.cls V1.6 and later.
% (Unless specifically asked to do so by the journal or conference you plan
% to submit to, of course. )


% *** CODE LISTING packages using Consola Font ***
\usepackage{listings}
\usepackage{inconsolata}
\lstset{
  basicstyle=\ttfamily,
}
\captionsetup[lstlisting]{font={small}, labelfont={bf}}

% *** FOOTNOTE package ***
\usepackage[bottom]{footmisc}
\setlength{\footnotemargin}{0.3em}


% correct bad hyphenation here
\hyphenation{op-tical networks semi-conduc-tor}
\hyphenation{PASCAL}

%***MARGIN***
% Do not set this margin
\usepackage[top=3cm,bottom=3cm,left=3cm,right=3cm]{geometry}

% *** HEADER & PAGE NUMBERING ***
% Do not set this header and page numbering.
\usepackage{fancyhdr}
\setlength{\headheight}{22pt}
\pagestyle{fancy}
\fancyhf{}
\setcounter{page}{1}
\fancypagestyle{firststyle}
{
   \fancyhf{}
   \chead{Jurnal Ilmu Komputer dan Informasi (Journal of Computer Science and Information) \\ 15/1 (2022), xx-xx. DOI: \url{http://dx.doi.org/10.21609/jiki.v15i1.xxx}}
   \fancyfoot[C]{\thepage}
}

\renewcommand{\headrulewidth}{0pt}
\fancyhead[LE]{\thepage \textbf{\space \textit{Jurnal Ilmu Komputer dan Informasi (Journal of Computer Science and Information),}} \textit{volume 15,} \\ \textit{ \textrm{ } issue 1, February 2022}}
\fancyhead[LO]{\textit{Author et.al. (left it blank for first submission), One Line Article Title}}
\fancyhead[RO]{\thepage}

%***Hyperref and URL line breaks***
\usepackage{hyperref}
\renewcommand{\UrlBreaks}{\do\-}

% Add horizontal line above footnotes suppressed by IEEE template
\makeatletter
\def\footnoterule{\kern-3\p@
  \hrule \@width 2in \kern 2.6\p@} % the \hrule is .4pt high
\makeatother
% Additional setup for figure captions
\def\figurename{Fig.}%
% \captionsetup[figure]{labelfont={bf},font=small}
\def\tablename{Table}%
% \captionsetup[table]{justification=justified,font=small, labelfont={bf}}


\begin{document}


%
% paper title
% Titles are generally capitalized except for words such as a, an, and, as,
% at, but, by, for, in, nor, of, on, or, the, to and up, which are usually
% not capitalized unless they are the first or last word of the title.
% Linebreaks \\ can be used within to get better formatting as desired.
% Do not put math or special symbols in the title.
\title{Simulation of Bandwidth-Efficient Federated Learning Architectures for Resource-Constrained Agricultural Networks in Indonesia}
% author names and affiliations
% use a multiple column layout for up to three different
% affiliations
\author{\IEEEauthorblockN{First Author\textsuperscript{1}, Second Author\textsuperscript{2}, Third Author\textsuperscript{1} (\textit{left it blank for first submission})\\\\}
	\IEEEauthorblockA{
		\normalfont \textsuperscript{1} Department, Faculty, University, Address, City, Zip Code, Country\\
		\textsuperscript{2} Research Group, Institution, Address, City, Zip Code, Country\\
		(\textit{left it blank for first submission})\\\\
		\textit{Email:{author@address.com} (\textit{left it blank for first submission})} 
	} \\
}

% conference papers do not typically use \thanks and this command
% is locked out in conference mode. If really needed, such as for
% the acknowledgment of grants, issue a \IEEEoverridecommandlockouts
% after \documentclass

% for over three affiliations, or if they all won't fit within the width
% of the page, use this alternative format:
% 
%\author{\IEEEauthorblockN{Michael Shell\IEEEauthorrefmark{1},
%Homer Simpson\IEEEauthorrefmark{2},
%James Kirk\IEEEauthorrefmark{3}, 
%Montgomery Scott\IEEEauthorrefmark{3} and
%Eldon Tyrell\IEEEauthorrefmark{4}}
%\IEEEauthorblockA{\IEEEauthorrefmark{1}School of Electrical and Computer Engineering\\
%Georgia Institute of Technology,
%Atlanta, Georgia 30332--0250\\ Email: see http://www.michaelshell.org/contact.html}
%\IEEEauthorblockA{\IEEEauthorrefmark{2}Twentieth Century Fox, Springfield, USA\\
%Email: homer@thesimpsons.com}
%\IEEEauthorblockA{\IEEEauthorrefmark{3}Starfleet Academy, San Francisco, California 96678-2391\\
%Telephone: (800) 555--1212, Fax: (888) 555--1212}
%\IEEEauthorblockA{\IEEEauthorrefmark{4}Tyrell Inc., 123 Replicant Street, Los Angeles, California 90210--4321}}

% use for special paper notices
%\IEEEspecialpapernotice{(Invited Paper)}
% make the title area
%\maketitle
\twocolumn[
{\csname @twocolumnfalse\endcsname \maketitle}
{\csname @twocolumnfalse\endcsname 
	%\input{abstract}
	\renewcommand{\abstractname}{Abstract}
	\begin{abstract}
		\noindent
		\normalfont 
		The deployment of federated learning (FL) in rural agricultural areas is often hindered by limited internet bandwidth and resource-constrained edge devices. This research aims to make collaborative AI more accessible by evaluating a communication-efficient FL architecture for weed and plant disease classification. The study utilizes Quantized Low-Rank Adaptation (QLoRA) to compress model updates and compares three training approaches: Centralized, Federated Averaging (FedAvg), and Distributed Low-Communication (DiLoCo). To simulate realistic rural constraints, the models are evaluated under an "Extreme Non-IID" data distribution, where isolated nodes possess entirely disjoint classes. The empirical results indicate that the DiLoCo approach significantly optimizes resource usage, achieving a 90% reduction in synchronization events and data transmission compared to FedAvg. Notably, despite communicating 10 times less, DiLoCo exhibits greater robustness to client drift, outperforming FedAvg in validation accuracy under extreme data heterogeneity. Furthermore, the 4-bit quantization ensures the training process requires a peak of only 748 MB of VRAM, proving its viability for deployment on low-cost agricultural edge devices. In conclusion, combining DiLoCo with QLoRA provides a highly robust, memory-efficient, and bandwidth-saving alternative for AIoT implementation in poorly connected agricultural environments. 
		\\\\
		\noindent
		\textbf{Keywords}: \textit{Federated Learning, Edge Computing, QLoRA, DiLoCo, Non-IID} \\\\
	\end{abstract}
	
	\renewcommand{\abstractname}{Abstrak}
% 	\begin{abstract}
% 		\noindent
% 		\normalfont 
% 		Abstrak dalam Bahasa Indonesia. Ditulis dengan font Times New Roman size 10 dan single spacing. Abstrak harus merangkum isi makalah, termasuk tujuan penelitian, metode penelitian, dan hasil, dan kesimpulan dari makalah. Abstrak tidak mengandung referensi dan/atau persamaan.Tidak boleh lebih dari 200 kata.
% 		\\\\
% 		\noindent
% 		\textbf{Kata Kunci}: \textit{terdiri dari 5 kata kunci} \\\\
% 	\end{abstract}
	
	%\begin{IEEEkeywords}
	%%IEEEtran, journal, \LaTeX, paper, template.
	%\end{IEEEkeywords}
}
%\vspace{1cm}
]

% For peer review papers, you can put extra information on the cover
% page as needed:
% \ifCLASSOPTIONpeerreview
% \begin{center} \bfseries EDICS Category: 3-BBND \end{center}
% \fi
%
% For peerreview papers, this IEEEtran command inserts a page break and
% creates the second title. It will be ignored for other modes.
\IEEEpeerreviewmaketitle
\thispagestyle{firststyle}
\graphicspath{{pics/}{../../../experiments/results/plots/}}

\section{Introduction}
% TODO: This section requires literature review and citations.
% Placeholder text below should be expanded with proper references.

The adoption of Artificial Intelligence (AI) in precision agriculture holds significant potential for improving crop management, particularly in weed detection and plant disease classification. However, deploying AI models in rural agricultural areas of Indonesia faces two critical barriers: limited internet bandwidth and resource-constrained edge devices. Traditional centralized machine learning approaches require all training data to be aggregated at a single server, raising concerns about data privacy, data sovereignty, and the impracticality of data transmission in bandwidth-scarce environments.

Federated Learning (FL) offers a promising paradigm where multiple distributed devices collaboratively train a shared model without exchanging raw data. Instead, only model updates (gradients or parameters) are communicated to a central server for aggregation. Despite this advantage, standard FL protocols such as Federated Averaging (FedAvg) still require frequent synchronization events that may be prohibitive in poorly connected networks.

This research investigates the feasibility of bandwidth-efficient federated learning for agricultural image classification by combining two key techniques: Quantized Low-Rank Adaptation (QLoRA) for model compression and Distributed Low-Communication (DiLoCo) for reduced synchronization frequency. The study evaluates a 3-way comparison of Centralized, FedAvg, and DiLoCo training approaches on a combined weed and plant disease dataset under an Extreme Non-IID data distribution that simulates realistic rural conditions where geographically isolated farming communities encounter entirely different agricultural problems.

\section{Methodology}

This section describes the experimental design used to evaluate bandwidth-efficient federated learning architectures for agricultural image classification. All experiments were conducted on an NVIDIA GeForce RTX 5050 Laptop GPU (8 GB VRAM) with CUDA 12.8 and PyTorch 2.10.0+cu128.

The experiment involves a controlled 3-way comparison of training approaches using QLoRA adapters on a combined multi-domain agricultural dataset with Extreme Non-IID label-based data partitioning. The three training approaches are Centralized (upper-bound reference), FedAvg (frequent-sync federated), and DiLoCo (rare-sync federated). The detailed configurations are summarized in Table~\ref{tab:experiment_config}.

\begin{table}[ht]
\centering
\caption{Experiment Configuration Summary. \medskip}
\label{tab:experiment_config}
\small
\begin{tabular}{ l  c  c  c  c }
\hline
Configuration & Farmers & Steps/Round & Rounds & Sync Events \\
\hline
Centralized + QLoRA & 1 & 10 epochs & --- & 0 \\
FedAvg + QLoRA & 10 & 50 & 200 & 200 \\
DiLoCo + QLoRA & 10 & 500 & 20 & 20 \\
\hline
\end{tabular}
\end{table}

The combined dataset consists of two agricultural image collections: WeedsGalore (3 weed classes, 156 images from 5-channel multispectral sources with RGB extraction) and PlantVillage (38 crop-disease pair classes, approximately 54,293 standard RGB leaf images). Label remapping assigns WeedsGalore classes 0--2 unchanged and PlantVillage classes 0--37 are offset to 3--40, yielding a total of 41 classes with approximately 43,516 images. The dataset is partitioned into approximately 34,812 training, 4,352 validation, and 4,352 test samples. All images are resized to 224$\times$224 pixels. The combined dataset is implemented via \texttt{ConcatDataset} with a \texttt{LabelOffsetDataset} wrapper for PlantVillage.

To simulate a realistic scenario where geographically isolated farmer communities encounter entirely different agricultural problems, we employ an Extreme Non-IID data distribution with disjoint label partitioning. Farmer Group A (Farmers 0--4) exclusively receives WeedsGalore data (classes 0--2, approximately 16--17 samples per farmer), while Farmer Group B (Farmers 5--9) exclusively receives PlantVillage data (classes 3--40, approximately 8,686--8,687 samples per farmer). This creates a worst-case scenario with zero class overlap between farmer groups and a massive sample imbalance ($\sim$17 vs $\sim$8,687 samples per farmer). The aggregated global model must therefore generalize to both domains despite each farmer learning from only one.

The base model is YOLOv11 Small Classification (YOLOv11s-cls), pretrained on ImageNet, with a total of approximately 9,449,000 parameters. The architecture comprises a YOLO11s backbone (layers 0--9 with Conv, C3k2, and C2PSA blocks, approximately 9,366,000 parameters), a classification convolutional layer Conv(256, 1280) with approximately 327,680 parameters, AdaptiveAvgPool2d(1), Flatten with Dropout(0.2), and a final Linear(1280, 41) classification head with 52,521 parameters.

The QLoRA adapter applies 4-bit NF4 quantization via bitsandbytes with float16 compute dtype. LoRA is configured with rank $r=4$, $\alpha=32$, dropout 0.05, targeting the \texttt{classifier.2} layer. This yields approximately 8,460 trainable adapter parameters, representing only $\sim$0.1\% of the total model parameters. Consequently, each synchronization event transmits only the adapter parameters rather than the full model, producing a per-sync bandwidth savings of 99.90\%.

\section{Results and Analysis}

This section presents the experimental results of the 3-way comparison across training loss convergence, validation performance, communication efficiency, GPU memory usage, and estimated economic impact.

Table~\ref{tab:training_loss} summarizes the training loss convergence for all three approaches. The Centralized + QLoRA approach shows steady monotonic loss reduction from 1.7545 to 1.1836 over 10 epochs. With access to all 43,000+ images and all 41 classes simultaneously, it converges smoothly without data distribution challenges. DiLoCo + QLoRA achieves the lowest training loss early (0.7657 at round 5) because 500 local steps allow deep fitting on each farmer's local data. However, severe client drift causes the training loss to rise to 4.6207 by round 20, which is expected under Extreme Non-IID since weed-farmers and disease-farmers pull the model in opposing directions. FedAvg + QLoRA initially decreases to 1.9879 (round 6) then diverges to 5.6326 by round 200. The frequent synchronization every 50 steps paradoxically worsens drift: constant averaging of weed-trained and disease-trained adapters creates conflicting gradient directions that prevent convergence.

\begin{table}[ht]
\centering
\caption{Training Loss Convergence Comparison. \medskip}
\label{tab:training_loss}
\small
\begin{tabular}{ l  c  c  c }
\hline
Metric & Centralized & FedAvg & DiLoCo \\
\hline
Initial train loss & 1.7545 & 3.3904 & 1.7801 \\
Final train loss & \textbf{1.1836} & 5.6326 & 4.6207 \\
Best train loss & \textbf{1.1727} & 1.9879 & 0.7657 \\
Trajectory & Monotonic $\downarrow$ & $\downarrow$ then $\uparrow$ & Deep fit then $\uparrow$ \\
\hline
\end{tabular}
\end{table}

Table~\ref{tab:val_performance} presents the validation performance. Centralized training dominates with a best validation accuracy of 71.39\% and macro F1-score of 0.6585, significantly outperforming both federated approaches. This is the expected result: centralized training sees all 41 classes, while federated farmers only ever see a subset. DiLoCo slightly outperforms FedAvg (48.81\% vs 44.96\% best validation accuracy; F1 0.3826 vs 0.3806) despite using 10 times fewer synchronization events. DiLoCo's rare synchronization allows each farmer to learn its local domain more deeply before averaging, reducing the destructive interference between weed-learned and disease-learned parameters. Both federated approaches plateau around 44--49\% accuracy, which is remarkable given the Extreme Non-IID constraint where random chance on 41 classes would be only 2.4\%.

\begin{table}[ht]
\centering
\caption{Validation and Test Performance Comparison. \medskip}
\label{tab:val_performance}
\small
\begin{tabular}{ l  c  c  c }
\hline
Metric & Centralized & FedAvg & DiLoCo \\
\hline
Best val accuracy & \textbf{0.7139} & 0.4496 & 0.4881 \\
Best val F1 (macro) & \textbf{0.6585} & 0.3806 & 0.3826 \\
Best val loss & \textbf{0.8270} & 2.9443 & 2.7475 \\
Test accuracy & \textbf{0.7082} & N/A & N/A \\
Test F1 (macro) & \textbf{0.6452} & N/A & N/A \\
\hline
\end{tabular}
\end{table}

Figure~\ref{fig:3way_results} illustrates the training loss convergence, validation accuracy progression, and bandwidth savings across the three approaches. The divergence of FedAvg loss after round 50 and the stability of DiLoCo despite rare synchronization are clearly visible.

\begin{figure*}[!t]
\centering
\includegraphics[width=\textwidth]{paper_figure_3way.png}
\caption{TaniFi QLoRA Federated Learning Comparison. (a) Training loss convergence showing Centralized monotonic decrease, FedAvg divergence after round 50, and DiLoCo deep fit followed by client drift. (b) Validation accuracy progression demonstrating Centralized dominance and DiLoCo's slight edge over FedAvg. (c) Bandwidth savings showing 99.9\% reduction for both federated approaches compared to full model synchronization.}
\label{fig:3way_results}
\end{figure*}

Table~\ref{tab:comm_efficiency} presents the communication efficiency analysis. DiLoCo achieves a 90\% reduction in total data transmitted compared to FedAvg (12.9~MB vs 128.9~MB) by using 10 times fewer synchronization events while maintaining the same per-sync adapter size of approximately 33.0~KB. Compared to traditional FedAvg with full YOLOv11s model synchronization (which would require 140.4~GB for 10 farmers over 200 rounds), the DiLoCo + QLoRA combination achieves a 99.99\% bandwidth reduction.

\begin{table}[ht]
\centering
\caption{Communication Efficiency Analysis. \medskip}
\label{tab:comm_efficiency}
\small
\begin{tabular}{ l  c  c  c }
\hline
Metric & FedAvg & DiLoCo & Reduction \\
\hline
Sync events & 200 & 20 & 10$\times$ fewer \\
Params per sync & $\sim$8,460 & $\sim$8,460 & Same \\
Per-sync data & $\sim$33.0 KB & $\sim$33.0 KB & Same \\
\textbf{Total data} & \textbf{128.9 MB} & \textbf{12.9 MB} & \textbf{90\% less} \\
BW saved vs full model & 99.90\% & 99.99\% & --- \\
\hline
\end{tabular}
\end{table}

Table~\ref{tab:vram} shows the GPU VRAM usage for each approach. Centralized training uses only 325.4~MB as it loads a single QLoRA model on the GPU. Federated approaches use 748.0~MB, which includes multi-farmer overhead from 10 farmers running sequentially on the same GPU. Both remain well within a 2,048~MB (2~GB) device constraint assumed for resource-constrained agricultural devices: the federated peak of 748~MB represents only 36.5\% of this limit. The 4-bit NF4 quantization keeps the backbone frozen and compact, confirming the viability of QLoRA with YOLOv11s-cls on devices with 1--2~GB of GPU memory.

\begin{table}[ht]
\centering
\caption{GPU VRAM Usage Comparison. \medskip}
\label{tab:vram}
\small
\begin{tabular}{ l  c  c  c }
\hline
Metric & Centralized & FedAvg & DiLoCo \\
\hline
Average VRAM (MB) & 325.4 & 748.0 & 748.0 \\
Peak VRAM (MB) & 325.4 & 748.0 & 748.0 \\
\hline
\end{tabular}
\end{table}

An economic impact estimate was calculated using average Indonesian rural mobile data costs of IDR 100/MB ($\sim$\$0.0067/MB) for a hypothetical network of 500 farmers conducting 10 training cycles annually. Traditional FedAvg with full YOLOv11s synchronization would cost approximately IDR 3,610,000,000 (\$240,667) per year. FedAvg + QLoRA reduces this to IDR 3,300,000 (\$220), and DiLoCo + QLoRA further reduces the annual cost to IDR 330,000 (\$22), representing a 99.99\% cost reduction from the traditional approach. This dramatic reduction makes collaborative AI training economically viable for smallholder farming communities with limited financial resources.

Several key findings emerge from the analysis. First, the Extreme Non-IID distribution severely challenges federated learning: both federated approaches achieve only 44--49\% validation accuracy compared to centralized's 71\%. This occurs because weed-farmers optimize exclusively for 3 weed classes while disease-farmers optimize for 38 disease classes, and model averaging creates a compromise model that is mediocre at both domains. Second, DiLoCo outperforms FedAvg under Extreme Non-IID conditions (+3.85 percentage points in best validation accuracy) while using 10 times fewer synchronization events. DiLoCo's rare synchronization preserves more domain-specific knowledge before averaging, whereas FedAvg's frequent averaging constantly undoes each farmer's local specialization. Third, despite the increase from 3 to 41 output classes, the QLoRA adapter size remains stable at approximately 8,460 trainable parameters because LoRA targets the \texttt{classifier.2} layer specifically, demonstrating QLoRA's scalability for multi-class agricultural tasks.

\section{Conclusion}

This study evaluated the feasibility of bandwidth-efficient federated learning for agricultural image classification on resource-constrained devices. Through a controlled 3-way comparison of Centralized, FedAvg, and DiLoCo training approaches using QLoRA adapters on a combined 41-class weed and plant disease dataset under Extreme Non-IID conditions, several important conclusions can be drawn.

The DiLoCo + QLoRA combination provides the most resource-efficient federated approach, achieving a 90\% reduction in total data transmission compared to FedAvg (12.9~MB vs 128.9~MB) while simultaneously achieving higher validation accuracy (48.81\% vs 44.96\%). Compared to traditional federated learning with full model synchronization, DiLoCo + QLoRA achieves a 99.99\% bandwidth reduction and reduces estimated annual mobile data costs from \$240,667 to \$22 for a 500-farmer network. The 4-bit NF4 quantization ensures a peak VRAM usage of only 748~MB, confirming viability for deployment on low-cost agricultural edge devices with 1--2~GB of GPU memory.

The Extreme Non-IID evaluation reveals a fundamental limitation of current federated aggregation methods when farmer data distributions are completely disjoint. The centralized upper bound of 71.39\% accuracy demonstrates the performance gap that future work should aim to close. Promising directions include implementing the DiLoCo outer optimizer with Nesterov momentum to mitigate client drift, testing intermediate Non-IID distributions to identify the tipping point where federation becomes beneficial, and applying per-domain evaluation to better understand domain-specific performance of the aggregated model.

In conclusion, combining DiLoCo with QLoRA provides a highly robust, memory-efficient, and bandwidth-saving alternative for AIoT implementation in poorly connected agricultural environments, making collaborative AI training accessible to smallholder farming communities in rural Indonesia.

\section*{Acknowledgement}
% Put your acknowledgement to people, funding, etc. just before the list of references, but only for the camera ready version.

\balance

\bibliographystyle{IEEEtran}
%\nocite{}
\bibliography{IEEEabrv,IEEEexample}

\end{document}
